% LIFECYCLE.tex — Chapter 3: the concept-agent lifecycle contract read against the factorization
% of Chapter 2 (FACTORIZATION.tex).
%
% Build standalone:  pdflatex LIFECYCLE.tex   (twice, for cross-references)
%
\documentclass[11pt,a4paper]{report}
\usepackage[T1]{fontenc}
\usepackage{lmodern}
\usepackage[margin=2.6cm]{geometry}
\usepackage{amsmath,amssymb,amsthm}
\usepackage{empheq}
\usepackage{pifont}
\usepackage{array}
\usepackage{booktabs}
\usepackage{enumitem}
\usepackage[hidelinks]{hyperref}
\usepackage{microtype}

\newcommand{\given}{\,|\,}
\newcommand{\R}{\mathbb{R}}

\begin{document}

\setcounter{chapter}{2}
\chapter{The lifecycle contract}

\begin{flushleft}
\textit{No\'e Zapata}
\end{flushleft}

\vspace{4pt}

\section{Introduction}

Every \texttt{<obj>\_concept} agent must satisfy a normative behavioural contract
(\texttt{CONCEPT\_\allowbreak AGENT\_\allowbreak LIFECYCLE.md}): what an agent may do to the shared
scene graph, and on
what evidence. Its five stages are \textsc{create}, \textsc{update}, \textsc{remove},
\textsc{history} and \textsc{ownership}.

That contract was written from experience rather than from theory. It exists because the same
defect class was found and fixed three separate times in three different agents, and because a
threshold that deleted objects the sensor provably could not see survived inside a hand-rolled
copy of a stage that already had a shared, correct implementation. Its purpose is that a stage
stated once is a stage a new agent cannot forget.

The previous chapter arrived at the same system from the opposite direction, by writing down the
joint distribution the agents approximate and the seven assumptions that make it computable. The
two descriptions were produced independently, which makes their comparison worth something: where
they agree, the contract has a derivation; where they disagree, one of them is missing a piece.

They agree in four stages out of five. The fifth, \textsc{create}, has no counterpart in the model
at all --- and neither does one clause of \textsc{update}. Those two gaps are the subject of
\S\ref{sec:create-gap} and \S\ref{sec:admissibility}, and they are, not coincidentally, where the
system's remaining hard thresholds are concentrated.

\section{The four stages that restate the factorization}

Writing $\theta_i$ for an instance's geometry, $e_i$ for its existence, $M$ for the per-cycle
mask-to-instance matching, and A1--A7 for the assumptions of Chapter~2:

\begin{center}
\begin{tabular}{@{}l l p{7.5cm}@{}}
\toprule
Stage & In the model & \\
\midrule
\textsc{update} & likelihood, A2, A5, A7 &
  ``a frame must earn the right to move geometry'' is the admissibility side of the measurement
  likelihood; the module the contract mandates is the recursive-Laplace engine itself \\[3pt]
\textsc{remove} & $e_i$, A6 &
  ``evidence, never a miss-counter'' requires removal to be driven by the existence factor, and A6
  is what makes that factor separable from geometry in the first place \\[3pt]
\textsc{history} & A7 &
  ``no constant may stand in for experience'' refines the Markov transition: it forbids a constant
  process noise $Q$ and requires an inferred, covariate-dependent one \\[3pt]
\textsc{ownership} & A4 &
  one agent owns the instances of one class, which is well defined only because A4 made the
  classes independent \\
\bottomrule
\end{tabular}
\end{center}

\noindent
\textsc{history} is the clearest case of the contract being ahead of a naive reading of the model.
A7 as stated permits any transition $P(\theta_i^t\given\theta_i^{t-1})$; the contract observes that
a \emph{constant} $Q$ applied every frame, observed or not, inflates covariance on the agent's
clock rather than on evidence, and requires the covariance growth to depend on what actually
happened. That is a strictly better model, arrived at from the symptom.

\section{\textsc{create} has no factor}
\label{sec:create-gap}

Birth adds an index to $\mathcal{I}_c$, and therefore adds a term to the product
$\prod_{i\in\mathcal{I}_c}$ in the assembled posterior of Chapter~2. Nothing in that posterior
decides that the term should be added.

\subsection{What the question actually is}

Deciding whether an object exists is a comparison between two models of different order: the world
with $n$ instances of this class, and the world with $n+1$. Under the variational treatment of
Chapter~2 that comparison has a definite form. Writing $F^{(n)}$ for the total free energy of the
class under $n$ instances,
\begin{align}
\Delta F \;=\; F^{(n+1)}-F^{(n)}
\;=\;&
\underbrace{\mathbb{E}_{q}\big[-\ln P(\bar Z\given\theta_{1:n+1})\big]
          -\mathbb{E}_{q}\big[-\ln P(\bar Z\given\theta_{1:n})\big]}_{\text{accuracy gained by explaining those points}}
\nonumber\\[4pt]
&+\;
\underbrace{\mathrm{KL}\!\left[\,q(\theta_{n+1})\,\|\,P(\theta_{n+1}\given e)\,\right]}_{\text{complexity paid for the new object}} .
\label{eq:deltaF}
\end{align}
The new instance is warranted exactly when $\Delta F<0$: when the measurements it claims are
explained better by a new object than by the clutter component and the existing ones, by more than
the cost of instantiating a new set of parameters. The second term is not a penalty bolted on for
regularisation --- it is the same complexity term that already appears in every per-instance free
energy, and it is what stops the model from explaining every stray mask with a new fridge.

\subsection{What is computed instead}

The contract requires that birth be evidence-accumulating rather than a counter, and mandates a
single shared implementation so it cannot be re-derived per agent. A detection matures a
\emph{pending candidate} over some number of frames, each frame contributing a weight. In the
refrigerator agent that weight is a product of several factors, each of which is a proxy for
something in \eqref{eq:deltaF}:

\begin{center}
\begin{tabular}{@{}l p{9.4cm}@{}}
\toprule
Factor & What it stands in for \\
\midrule
detector confidence & the likelihood's weight on the class label \\
range attenuation   & measurement precision, which falls with distance \\
frame admissibility & whether this frame may move geometry at all \\
shape plausibility  & the prior term $P(\theta_{n+1}\given e)$, evaluated on the raw cloud \\
room-envelope support & the same prior, evaluated against the wall polygon \\
residual surprise (optional) & unexplained occupancy under the detection: the accuracy term \\
\bottomrule
\end{tabular}
\end{center}

\noindent
The list is a reasonable decomposition. Every entry corresponds to a term that would appear in
$\Delta F$, and the shared implementation means all agents birth on the same terms.

\subsection{What the surrogate costs}

Three things are lost by accumulating a weighted count rather than evaluating \eqref{eq:deltaF}.

\paragraph{The units are gone.} A streak is a dimensionless number compared against a frame count.
$\Delta F$ is in nats, and can therefore be traded against any other quantity in the model --- the
existence log-odds, a merge decision, the cost of a viewpoint. Because the birth score is not in
nats, birth cannot participate in any of those trades, and the boundary it is compared against is
a threshold with no principled value.

\paragraph{Birth and death are not the same currency.} Removal is driven by an existence log-odds
accumulated from occupancy carve and silhouette evidence. Birth is driven by a streak. The two
decisions concern the same proposition --- does this object exist --- and are taken in
incommensurable units, so an object can be born on evidence that would not have kept it alive, and
the asymmetry is invisible to both mechanisms.

\paragraph{The comparison is available and not made.} $\Delta F$ is not out of reach. A candidate
that has matured has an associated set of measurements; instantiating a provisional instance,
fitting it, and reading off the resulting free energy against the null is the same computation the
agent already performs every cycle for its live instances. What is missing is the plumbing, not
the theory.

\subsection{Where this shows up}

The refrigerator agent carries a subsystem whose sole purpose is to undo births that should not
have happened: shape plausibility accumulated into the existence channel, plus a soft singleton
inhibition so a strongly-supported instance suppresses weaker duplicates. It works, and it is
carefully built. It is also a corrective loop that exists because the birth decision could not
weigh the evidence properly at the time it was taken --- and its own criterion is, again, in
different units from the one that created the object.

A birth gate that computed \eqref{eq:deltaF} on the candidate's accumulated measurements would
subsume the plausibility multiplier, the envelope support, the residual-surprise fusion and the
retraction subsystem, because all four are estimating parts of the same scalar.

\section{The \textsc{update} admissibility gate}
\label{sec:admissibility}

The second gap is smaller and of the same kind. The contract requires a frame to be
\emph{resolvable}, and the robot \emph{still} or the mask \emph{well-centred}, before that frame
may move the geometric mean.

In the model there is no such switch. A poor frame should carry a large observation covariance and
lose to the prior on its own, continuously and without a boundary. The gate exists because the
covariance model is not yet trusted to do that unaided: motion corruption, truncation and
off-centre bias are known to inflate the error in ways the current $R$ under-reports, so a hard
predicate stands in for the missing precision.

This one is closer to being closed than \textsc{create}. Each condition in the predicate names a
covariate that already has a measured proxy --- ego-motion magnitude, truncated fraction, centroid
radius --- and each is already carried on the mask. Folding them into $R$ instead of into an
\texttt{if} is a change of the same shape as the one \textsc{history} made for $Q$: replace a
constant-or-switch with a covariate-dependent term and let the inference do the gating.

\section{Summary}

Four of the five lifecycle stages are assumptions of the factorization wearing operational
clothes, and one of them --- \textsc{history} --- improves on the assumption as stated. The two
remaining rules, the birth criterion and the update admissibility gate, are doing work the model
has no term for, and are correspondingly the places where the project's own preference for a
continuous covariance over a threshold is not yet satisfied.

Both gaps have the same remedy in outline: compute the quantity the heuristic is approximating.
For \textsc{update} that means moving three known covariates into the observation covariance. For
\textsc{create} it means evaluating a free-energy difference that the agent is already equipped to
compute, and letting birth, retraction and removal be decided in one currency instead of three.

\end{document}
